\documentclass[11pt,a4paper]{article}
\usepackage[utf8]{inputenc}
\usepackage{geometry}
\geometry{margin=1in}
\usepackage{hyperref}
\usepackage{graphicx}
\usepackage{booktabs}
\usepackage{tabularx}
\usepackage{float}
\usepackage{setspace}

\title{\textbf{DSM A2 Final Report: Yelp Open Dataset}}
\author{Shivansh Verma}
\date{\today}

\begin{document}
\maketitle

\tableofcontents
\newpage


\section{Data Acquisition and MongoDB Schema Design}

\subsection*{1. Core ETL Logic \& Data Loading}
The extraction, transformation, and loading (ETL) pipeline was developed in Python utilizing the \texttt{pymongo} library. Instead of a direct import mirroring the exact JSON constraints, we apply specific transformations defined in \texttt{etl/load\_yelp.py}. The pipeline uses generators to stream huge \texttt{.json} files, formatting attributes like \texttt{date} strings into proper MongoDB ISODate objects, restructuring nested values like \texttt{checkins} into formal Datetime Arrays, and merging native IDs to MongoDB's intrinsic \texttt{\_id} fields to inherently optimize index sizing as data loads.

The script runs sequentially via:
\begin{verbatim}
python3 etl/load_yelp.py --data-dir ../data
\end{verbatim}

\subsection*{2. Collection Structure and Inter-relations}
The schema was modeled strictly based on the relational and analytical needs of the provided dataset parameters. Instead of mimicking the six source JSON files 1:1, the design separates bounded attributes (embedded) from unbounded elements (referenced) to prevent hitting MongoDB's 16MB document size limit as relationships scale.

Below is an overview of the six primary collections, their respective roles, document identifiers, and referential structures.

\subsubsection*{Businesses Collection}
\textbf{Purpose:} Acts as the central anchor for the database, housing core attributes, categorizations, and static metadata for various commercial locations.\\
\textbf{Identifier Field:} \texttt{\_id} (String, mapped directly from Yelp's native \texttt{business\_id}).\\
\textbf{Linking Relationships:} Is the parent entity. Other collections (like reviews, tips, and checkins) reference a business by its \texttt{\_id}.\\
\textbf{Key Fields:} \texttt{name}, \texttt{city}, \texttt{stars}, \texttt{review\_count}, \texttt{is\_open}. Attributes like \texttt{categories} (Array) and \texttt{hours} (Object) are strictly bound to the business and do not continuously grow, making them ideal to embed directly for localized queries.

\subsubsection*{Users Collection}
\textbf{Purpose:} Stores user metadata, engagement metrics, and social network ties.\\
\textbf{Identifier Field:} \texttt{\_id} (String, mapped directly from Yelp's native \texttt{user\_id}).\\
\textbf{Linking Relationships:} The \texttt{friends} field contains an array of strings referencing other \texttt{\_id} fields within the same collection (N:M self-referential). It is also globally referenced by the Reviews and Tips collections.\\
\textbf{Key Fields:} \texttt{name}, \texttt{yelping\_since} (Date), \texttt{elite} (Array), \texttt{average\_stars}, and various metrics. Due to the scale of social connectivity, \texttt{friends} arrays are populated strictly via references rather than embedded documents to avoid extreme growth amplification.

\subsubsection*{Reviews Collection}
\textbf{Purpose:} Logs independent evaluations made by users on specific businesses.\\
\textbf{Identifier Field:} \texttt{\_id} (String, mapped from \texttt{review\_id}).\\
\textbf{Linking Relationships:} Contains \texttt{user\_id} and \texttt{business\_id} fields, which inherently reference the Users and Businesses collections (N:1 mapping to both).\\
\textbf{Key Fields:} \texttt{stars}, \texttt{date} (Date object), \texttt{text}, \texttt{useful}. Referencing is strictly necessary here. If reviews were embedded within a user or business, heavily trafficked locations could easily crash string limitations.

\subsubsection*{Checkins Collection}
\textbf{Purpose:} Tracks longitudinal visitation trends mapped over time.\\
\textbf{Identifier Field:} \texttt{\_id} (String, mapped directly from \texttt{business\_id}), forming a rigid 1:1 schema extension with the Business entity.\\
\textbf{Linking Relationships:} Because the \texttt{\_id} maps to \texttt{business\_id}, it explicitly references the Businesses collection.\\
\textbf{Key Fields:} \texttt{dates} (Array of Date objects). Although checkins are technically bounded to a business, separating them into a new collection prevents the core business documents from massive index bloat generated by extremely large, actively expanding sub-arrays of timestamps.

\subsubsection*{Tips Collection}
\textbf{Purpose:} Captures smaller, qualitative offhand recommendations distinct from formal reviews.\\
\textbf{Identifier Field:} \texttt{\_id} (Auto-generated assigned by MongoDB).\\
\textbf{Linking Relationships:} Utilizes distinct \texttt{user\_id} and \texttt{business\_id} fields to bridge across the Users and Businesses collections.\\
\textbf{Key Fields:} \texttt{text}, \texttt{date} (Date object), \texttt{compliment\_count}. Like reviews, tips are unconstrained and inherently must maintain a normalized cross-reference.

\subsubsection*{Photos Collection}
\textbf{Purpose:} Contains metadata and classifications for user-uploaded images associated with a business.\\
\textbf{Identifier Field:} \texttt{\_id} (String, mapped directly from \texttt{photo\_id}).\\
\textbf{Linking Relationships:} The \texttt{business\_id} field provides a N:1 referenced mapping directly to the Businesses collection.\\
\textbf{Key Fields:} \texttt{caption}, \texttt{label} (String classification). Standalone referencing prevents massive accumulation of string URLs or binary blobs within the base business document.

\subsection*{3. Schema Design Justifications (Embedding vs. Referencing)}

\subsubsection*{(a) Relationship Design Choices}
\begin{itemize}
    \item \textbf{Business $\rightarrow$ Attributes, Categories, Hours (Embedded):} These entities inherently map 1:1 to a specific business. Their data volume is strictly bounded (e.g., there are only 7 days in a week for hours, and a finite set of attributes). Embedding these properties prevents expensive cross-collection joins just to retrieve basic business profile information.
    \item \textbf{User $\rightarrow$ Friends (Referenced):} The Yelp social graph is massive and highly interconnected (N:M). Embedding a user's friends list as full documents would exponentially duplicate data. Even embedding an array of \texttt{user\_id} strings risks exceeding MongoDB's 16MB document size limit for highly active "elite" users. Thus, referencing is strictly required.
    \item \textbf{Business $\rightarrow$ Reviews \& Checkins (Referenced):} Yelp reviews and checkins grow unboundedly over time (1:N mapping). Embedding thousands of reviews directly inside a single \texttt{business} document would inevitably breach the 16MB limit, severely degrading read/write performance long before failing entirely. Keeping these in standalone collections referencing \texttt{business\_id} is the only scalable approach.
    \item \textbf{User $\rightarrow$ Reviews \& Tips (Referenced):} Similar to businesses, a highly active user can generate an unbounded number of reviews and tips over their tenure. These must uniquely exist in their own collections, maintaining references back to the contributing \texttt{user\_id}.
\end{itemize}

\subsubsection*{(b) Read and Write Trade-offs}
\textbf{Embedding Trade-offs (Business Attributes/Categories/Hours):}
\begin{itemize}
    \item \textbf{Read Advantage:} Yields unmatched read performance. A single query to the Businesses collection retrieves the entire business profile, navigating latency from \texttt{\$lookup} joins.
    \item \textbf{Write Disadvantage:} Granular updates (like changing a single attribute) require rewriting the entire embedded object mapping. However, since business static attributes rarely change, this write penalty is highly negligible.
\end{itemize}

\textbf{Referencing Trade-offs (Reviews/Checkins/Friends):}
\begin{itemize}
    \item \textbf{Write Advantage:} Excellent write performance and infinite scalability. When a new review is posted, it simply inserts a new document in the Reviews collection without ever locking or modifying the massive parent Business or User document.
    \item \textbf{Read Disadvantage:} Requires application-level joins or aggregation pipeline \texttt{\$lookups}. Fetching a business alongside all its reviews demands querying multiple collections, increasing memory usage and retrieval latency. This trade-off is mitigated by strategic indexing.
\end{itemize}

\subsubsection*{(c) Indexing Strategy for Section 4 Queries}
To efficiently execute the analytical queries mapped out in Section 4 of the assignment, the following targeted indexes were established inside the ETL loading script:

\begin{enumerate}
    \item \texttt{db.businesses.createIndex(\{ "city": 1, "categories": 1 \})}
    \begin{itemize}
        \item \textbf{Justification:} Section 4 Queries 1 \& 2 strictly demand isolating specific cities and business categories to compute safety and ratings. This compound index allows MongoDB to instantly filter down the 150,000+ businesses dataset without triggering an expensive full collection scan (COLLSCAN).
    \end{itemize}
    \item \texttt{db.reviews.createIndex(\{ "business\_id": 1, "date": 1 \})}
    \begin{itemize}
        \item \textbf{Justification:} Query 2 analyzes rating trends \textit{over time}. This compound index supports extremely fast sequential retrievals natively sorted by date for any given business, removing the need for an expensive in-memory sort stage during the aggregation pipeline.
    \end{itemize}
    \item \texttt{db.reviews.createIndex(\{ "user\_id": 1 \})}
    \begin{itemize}
        \item \textbf{Justification:} Queries 5 \& 6 require aggregating a user's tenure and elite status against their review behavior (average review length and usefulness ratio). This index enables rapid resolution of all reviews mapped to a specific user during a \texttt{\$lookup} join originating from the Users collection without scanning 7 million generic reviews.
    \end{itemize}
    \item \texttt{db.checkins.createIndex(\{ "\_id": 1 \})} \textit{(Auto-generated mapped from business\_id)}
    \begin{itemize}
        \item \textbf{Justification:} \texttt{checkins} use \texttt{business\_id} directly as their explicit \texttt{\_id}, granting them an automatic, uniquely optimal index at scale. For Query 7 (correlating check-in visit frequency with business star ratings), we natively rely on this \texttt{\_id} index to flawlessly execute a 1:1 optimal pipeline join with the \texttt{businesses} collection on the same key.
    \end{itemize}
\end{enumerate}

\end{document}
